\subsection{Recovery}
Recovery has to make sure that, even in the case of failures, the changes from transactions that have not completed are not visible
and those of committed transactions are recorded and visible.
We define undesirable situations and classify histories according to these situations.
Furthermore, we want to define how histories can avoid such situations.


\subsubsection{Recovery procedures}
\begin{enumerate}[label=R\arabic*]
    \item Undo all changes from transactions
    \item Redo the changes from committed transactions
    \item Undo the changes that remain in the system from active transactions
    \item Read consistent snapshot from backup, and, if available, apply log
\end{enumerate}
R1 is used for abort/rollback of single transactions, R2 and R3 are used on system crash (where we lose the memory, but disks survive)
and R4 is used for recovery from a system crash with disk loss.

This means that we can undo the changes of an active transaction (we keep a copy of the value before it was modified; Shadow pages is a mechanism that would support this),
we can redo changes made by committed transactions (we keep a persistent record of all the changes made by committed transactions, this why there is a redo log)
and we have a consistent snapshot to start with (either move forward by taking a snapshot and apply committed transactions,
or move backward by taking the current state and go back to a consistent state by undoing changes from transactions that should be there).


\subsubsection{Undo - Redo}
The changes of an aborted transaction are undone typically by restoring the \bi{before image} of the modified value.

Changes are redone by restoring the \bi{after image} of the modified value.

Databases typically log transactions by keeping before and after images of their changes.


\subsubsection{Recovery on histories}
\label{sec:recovery-of-histories}
\begin{examdetails}
    They like to ask about these, so be sure to understand this properly. Same goes for Conflict Serializability for histories (see \ref{sec:tm-history})
\end{examdetails}
A transaction $T_1$ reads from another transaction $T_2$ if $T_1$ reads a value written by $T_2$ at a time when $T_2$ was not aborted.

We have:
\begin{itemize}
    \item \bi{Recoverable} (RC) history: If $T_i$ reads from $T_j$ and commits, then $c_j < c_i$ (same goes for aborts)
    \item \bi{Avoids Cascading Aborts} (ACA) history: If $T_i$ reads $x$ from $T_j$, then $c_j < r_i[x]$
    \item \bi{Strict} (ST) history: If $T_i$ reads from, or overwrites, a value written by $T_j$ then $c_j < o_i[x]$ or $a_j < o_i[x]$, with $o = r$ or $o = w$ ($o$ is an operation)
\end{itemize}


\paragraph{Recoverable}
There is no need to undo a committed transaction because it read the wrong data and they commit in their serialization order.

Suppose we have the following sequence: $\ldots w_2[x] r_1[x] c_1 \ldots$.
If $T_2$ aborts, then $T_1$ has read invalid data and since $T_1$ has committed, the data has returned to the user.
Recovery now involves correcting the application that has processed the data and might have acted upon it, which typically is a very expensive procedure.
This is avoided by making sure that $T_1$ does not commit until $T_2$ commits and if $T_2$ aborts, then we can safely abort $T_1$,
since the results have not been returned to the user yet.


\paragraph{ACA}
Aborting a transaction does not cause aborting others and transactions only read from committed transactions

Suppose we have the following sequence: $\ldots w_2[x] r_1[x] a_2 \ldots$.
When $T_2$ aborts, we have to abort $T_1$ because it has read invalid data (the change made by $T_2$, which will be undone as part of rollback of $T_2$).
If those cascading aborts happen regularly, the performance will degrade and possibly further transactions will be aborted.
This is avoided by making sure that uncommitted data is never read (\textit{read committed}).


\paragraph{Strict}
Undoing a transaction does not undo the changes of other transactions. Transactions further don't read or overwrite updates of uncommitted transactions.

Suppose we have the following sequence: $\ldots w_2[x] w_1[x] a_2 \ldots$.
If we undo $w_2[x]$, we would also remove the changes made by $T_1$. Thus, undoing $T_2$ implies aborting $T_1$.
Like in ACA, non-strict execution leads to cascading aborts of transactions.
This is avoided by not letting a value be read or updated unless it is committed.


\paragraph{Why recoverability matters}
Fixing errors with recovery typically are very expensive:
\begin{itemize}
    \item Not RC: running my program issues a report, but the data that I got was later removed because a different program aborted $\rightarrow$ report is invalid
    \item Not ACA: Thrashing behaviour when transactions keep aborting each other (my program slows down, or makes no progress at all because of other programs)
    \item Not Strict: recovery after a failure becomes very complex (or impossible)
\end{itemize}
